\documentclass[11pt, oneside]{article}   	% use "amsart" instead of "article" for AMSLaTeX format
\usepackage{geometry}                		% See geometry.pdf to learn the layout options. There are lots.
\geometry{letterpaper}                   		% ... or a4paper or a5paper or ... 
%\geometry{landscape}                		% Activate for rotated page geometry
%\usepackage[parfill]{parskip}    		% Activate to begin paragraphs with an empty line rather than an indent
\usepackage{graphicx}				% Use pdf, png, jpg, or eps§ with pdflatex; use eps in DVI mode
								% TeX will automatically convert eps --> pdf in pdflatex		
\usepackage{amssymb}
\usepackage{amsmath}
\usepackage{listings}
\usepackage{color}
\usepackage{geometry}
\usepackage{booktabs}

\usepackage{array}
\usepackage{longtable}
\usepackage{xcolor}
\usepackage{colortbl}
\usepackage{caption}
\usepackage{hyperref}

\usepackage{minted}

\definecolor{dkgreen}{rgb}{0,0.6,0}
\definecolor{gray}{rgb}{0.5,0.5,0.5}
\definecolor{mauve}{rgb}{0.58,0,0.82}

\lstset{frame=tb,
  language=Java,
  aboveskip=3mm,
  belowskip=3mm,
  showstringspaces=false,
  columns=flexible,
  basicstyle={\small\ttfamily},
  numbers=none,
  numberstyle=\tiny\color{gray},
  keywordstyle=\color{blue},
  commentstyle=\color{dkgreen},
  stringstyle=\color{mauve},
  breaklines=true,
  breakatwhitespace=true,
  tabsize=3
}


\begin{document}

\centerline {\bf Thoughts on the analysis of the violin}
This document has an in-depth explanation of all the math and file formats I use in the ObieApp.
\section{Frequency Response Functions}
\subsection{What is a FRF?}
We want to express the violin mathematically, we can view it as a transfer function (also called the frequency response function)- it converts some input (in this case a force from the hammer) to some output (in this case sound or a pressure fluctuation):
\begin{equation}
f(t) -> p(t)
\end{equation}
or taking the Fourier Transform, we can find the transfer function for the violin as:
\begin{equation}
H(f) = {P(f)\over F(f)}
\end{equation}
where $P(f)$ and $F(f)$ are the Fourier transforms of pressure and force and have real and imaginary parts.  Remember that the Fourier Transform is defined as:
\begin{equation}
G(f) = \int_{-\infty}^{\infty} g(t) e^{-i2\pi ft}dt
\end{equation}
I can simplify - get the denominator to be real, by multiplying by $F^*$:
\begin{equation}
H(f) = {P(f)\over F(f)}={PF^*\over FF^*}= H1 = {PP^*\over FP^*} = H2
\end{equation}

\subsection{Least Squares Fit}
In reality, we take 5-10 hits per position so that we can get a ``best fit" reading.  We know that:
\begin{equation}
P(f) = H(f)F(f)
\end{equation}
So the error of the best fit would be looking at the measured values ($F_i \text{ and }P_i$):
\begin{equation}
\epsilon_i(f) = P_i - H(f)F_i
\end{equation}
Now realizing that $H(f)$ is a complex number ($a+ib$) and that $F,P$ are complex as well, gives (remember that the magnitude of the product equals the product of the magnitudes):
\begin{equation}
    \begin{split}
        \epsilon_i^2(f) &= (P_i - H(f)F_i)(P_i - H(f)F_i)^*\\
        &= |P_i|^2 -H(f)F_iP^*_i - P_iH^*(f)F^*_i + |H(f)F_i|^2 \\
        &= |P_i|^2 -(a+ib)F_iP^*_i - P_i(a-ib)F^*_i + |(a+ib)F_i|^2\\
        &= |P_i|^2 -(a+ib)F_iP^*_i - P_i(a-ib)F^*_i + (a^2+b^2)F_iF_i^*
    \end{split}
\end{equation}
Next we sum over all the readings and then find the minimum by differentiating with respect to $a$ and $b$ and setting them equal to 0 to get:
\begin{equation}
    \begin{split}
       \frac{\partial \sum \epsilon_i^2(f)}{\partial a} = 0= 0 -\sum F_iP^*_i - \sum P_iF^*_i + 2a\sum F_iF_i^*\\
       \frac{\partial \sum \epsilon_i^2(f)}{\partial b} = 0=0 -i\sum F_iP^*_i + i\sum P_iF^*_i  +2b\sum F_iF_i^*
    \end{split}
\end{equation}
multiplying the second one by $i$ and then adding you get:
\begin{equation}
    \begin{split}
       0=-\sum F_iP^*_i - \sum P_iF^*_i + 2a\sum F_iF_i^*\\
       0= \sum F_iP^*_i -\sum P_iF^*_i +2bi\sum F_iF_i^*
    \end{split}
\end{equation}
I can add them together and get:
\begin{equation}
    \begin{split}
       0 &= -2 \sum P_iF^*_i + 2a\sum F_iF_i^* +2bi\sum F_iF_i^*\\
        \sum P_iF^*_i &= H(f)\sum F_iF_i^*\\
    \end{split}
\end{equation}
or 
\begin{equation}
    \begin{split}
        H_1(f) = \frac{\sum P_iF^*_i}{\sum F_iF_i^*} \\
    \end{split}
\end{equation}
Likewise we could look for reducing the effect of noise on the input -- or using, as the error, 
\begin{equation}
\epsilon_i(f) = F_i - \frac{P_i}{H(f)}
\end{equation}
And do the analysis again and you will get:
\begin{equation}
    \begin{split}
        \epsilon_i^2(f) &= (F_i - \frac{P_i}{H(f)})(F_i - \frac{P_i}{H(f)})^*\\
        &= F_iF_i^* - F_i\frac{P_i^*}{H^*}- F^*_i\frac{P_i}{H}+\frac{P_i^*}{H^*}\frac{P_i}{H}\\
        &= F_iF_i^* - F_iP_i^*(a-ib)^{-1}- F^*_iP_i(a+ib)^{-1}+P_iP_i^*(a^2+b^2)^{-1}
    \end{split}
\end{equation}
Again summing and differentiating:
\begin{equation}
    \begin{split}
       \frac{\partial \sum \epsilon_i^2(f)}{\partial a} = 0= 0 +(a-ib)^{-2}\sum F_iP^*_i +(a+ib)^{-2}\sum P_iF^*_i -(a^2+b^2)^{-2}2a\sum P_iP_i^*\\
       \frac{\partial \sum \epsilon_i^2(f)}{\partial b} = 0= 0 -i(a-ib)^{-2}\sum F_iP^*_i +i(a+ib)^{-2}\sum P_iF^*_i -(a^2+b^2)^{-2}2b\sum P_iP_i^*\\
    \end{split}
\end{equation}
multiplying the second equation through by $i$ gives you:
\begin{equation}
    \begin{split}
       0= (a-ib)^{-2}\sum F_iP^*_i +(a+ib)^{-2}\sum P_iF^*_i -(a^2+b^2)^{-2}2a\sum P_iP_i^*\\
       0= (a-ib)^{-2}\sum F_iP^*_i -(a+ib)^{-2}\sum P_iF^*_i -(a^2+b^2)^{-2}2ib\sum P_iP_i^*\\
    \end{split}
\end{equation}
Adding them together and reshuffling:
\begin{equation}
    \begin{split}
       H(a^2+b^2)^{-2}\sum P_iP_i^*&= (a-ib)^{-2}\sum F_iP^*_i \\
       H\sum P_iP_i^*&= \frac{(a+ib)^2(a-ib)^2}{(a-ib)(a-ib)}\sum F_iP^*_i \\
       \sum P_iP_i^*&= H\sum F_iP^*_i \\
       H_2(f) &= \frac{\sum P_iP_i^*}{\sum F_iP^*_i} \\
    \end{split}
\end{equation}
gives us the best fit in the other direction (assuming that the input is noisy).  If the data is clean and the noise is truly random, these two values for $H$ should be very close and therefore their ratio, the coherence, should be 1.
\begin{equation}
C = \frac{H_1}{H_2}
\end{equation}
If it is not 1, then one (or both) of your signals is noisy (not a great least squares fit). George then used the geometric mean.
\begin{equation}
H\sim \sqrt{H_1 H_2}
\end{equation}

\subsection{Claude's explanation: Coherence}
Coherence is essentially a measure of how consistently a system behaves across repeated measurements.

Think of it this way: you tap a violin with a hammer and measure the input force and the output sound. If you tap it again with a slightly different force, does the ratio of output to input stay the same? If the system is well-behaved (linear, no noise), it should — and the coherence will be close to 1. If something is messy (nonlinearity, background noise, a loose part rattling), the ratio will vary between hits, and coherence drops toward 0.
More concretely, coherence answers: "how much of the output can be explained by the input?"
The formula from the document makes this intuitive:
\begin{equation}
C = \frac{H1}{H2}=\frac{|\sum P_i F_i^*|^2}{\sum F_i F_i^*  \sum P_i P_i^*}
\end{equation}
The numerator is the cross-spectrum squared — how much F and P move together. The denominator is the product of the individual power spectra — how much each signal varies on its own. It's essentially a frequency-domain analog of $R^2$ from linear regression.
A few practical rules of thumb:
\begin{itemize}
\item C $\sim$ 1: the output is well explained by the input — clean measurement, linear system
\item C $\sim$ 0: the output has little to do with the input at that frequency — probably noise, or the system simply doesn't respond there
\item C between 0 and 1: partial explanation — could be noise, nonlinearity, or multiple inputs contributing to the output
\end{itemize}
One important subtlety: coherence is frequency-dependent. A violin might have C $\sim$ 1 near a resonance (where it responds strongly and predictably) but low coherence in a frequency range where it barely responds and the output is dominated by noise. So you'd typically plot it across the whole frequency range rather than summarize it as a single number.

\section{Band Averages}
One can look at an average distribution of the energy by breaking the FRF up into bands and averaging the FRF magnitudes over that area.
\begin{equation}
    \begin{split}
        \bar b = \frac{1}{N}\sum_{f_{lo}}^{f_{hi}} \sqrt{H(f)H^*(f)} \\
        \bar b_{dB} = 20 log(\frac{\bar b}{1})
    \end{split}
\end{equation}
We also calculate the centroid of the band - that is the first moment:
\begin{equation}
c(f) = \frac{\sum_{f_{lo}}^{f_{hi}} f\sqrt{HH^*}}{\sum_{f_{lo}}^{f_{hi}} \sqrt{HH^*}}
\end{equation}
which tells you something about the variation of the energy within the band.
or in Python:
\begin{minted}{python}
    f_lo  = float(band['start'])
    f_hi  = float(band['end'])
    mask = (freq >= f_lo) & (freq <= f_hi)
    H_band   = H_mag[mask]  #NOT in dB
    f_band   = freq[mask]
    avg_db   = 20.0 * math.log10(max(float(np.mean(H_band)), 1e-12))
    centroid = float(np.sum(f_band * H_band) / np.sum(H_band))
\end{minted}

\section{Convolutions}
One can estimate what the instrument will sound like by convolving the FRF with a recording of a song - but using a force sensor at the bridge rather than a microphone.  So if you want to convolve two functions $f(t)$ and $g(t)$
\begin{equation}
    y(t) = \int_{-\infty}^\infty f(\tau)g(t-\tau)d\tau
\end{equation}

\subsection{What they are - from Claude}

\textbf{What convolution does: }
A violin has resonances — certain frequencies ring longer and louder than others. If you record the violin's frequency response (the FRF), you have a fingerprint of exactly how it shapes sound.
Convolution applies that fingerprint to any audio signal. Feed in a Tchaikovsky recording; get out what it would sound like played through or colored by that violin's resonance structure.

\textbf{How it works: }
Sound is made of many frequencies added together. The FRF tells you how much the violin boosts or cuts each frequency — and by how much it time-shifts it (phase part).

Note that you can speed up the calculation by multiplying the two waveforms in frequency space and then inverse Fourier transforming them into time space:

\begin{equation}
    \begin{split}
        y(t) &= iFFT(FFT(g(t))\times FFT(f(t))\\
        &= \int_{-\infty}^\infty [\int_{-\infty}^\infty g(\tau)e^{-i2\pi f\tau}d\tau][\int_{-\infty}^\infty f(\tau)e^{-i2\pi f\tau} d\tau]e^{i2\pi ft}df
    \end{split}
\end{equation}
or in Python:
\begin{minted}{python}
    y = np.fft.irfft(
        np.fft.rfft(wav, n_fft) * np.fft.rfft(ir.astype(np.float64), n_fft),
        n_fft,)[:N].real
\end{minted}

\textbf{What you hear: }
The output is the input signal with the violin's character stamped onto it — its resonant peaks amplified, its dead spots suppressed, its natural decay times imposed on every note. If the FRF came from a particularly resonant instrument, the result sounds richer; from a dead one, duller.

\subsection{minimum phase reconstruction from Claude}
In the case that you do not have the complex FRF, you can still estimate the convolution by realizing for a minimum phase system, the log-magnitude and phase are a Hilbert transform pair. So if you know one, you can compute the other.  We do not need to make use of this since everything is stored as complex FRFs.

\section{File Formats}

\subsection{TRF files}
Parse and build binary \texttt{.trf} files (Rational Acoustics / MtxVec format).

\subsubsection*{Header Format}
Little-endian, 110 bytes total:

\begin{tabular}{ll}
\toprule
\textbf{Type} & \textbf{Fields} \\
\midrule
\texttt{I}    & \texttt{index} \\
\texttt{4d}   & \texttt{Xr, Yr, Xactual, YActual} \\
\texttt{2s}   & \texttt{char\_str} \\
\texttt{3d}   & \texttt{Hz\_Resolution, Start\_Freq, End\_Freq} \\
\texttt{11f}  & \texttt{fComplex, fLength, \ldots} (9 unused floats) \\
\texttt{4s}   & \texttt{caption} \\
\bottomrule
\end{tabular}

\subsubsection*{Data (from byte 110)}
\begin{itemize}
    \item \textbf{Complex:} alternating \texttt{float64} real/imag pairs
    \item \textbf{Real:} plain \texttt{float64} values
\end{itemize}


\subsection{File Formats}
\definecolor{headerblue}{RGB}{30, 80, 140}
\definecolor{sectiongray}{RGB}{220, 225, 230}
\definecolor{rowlight}{RGB}{240, 244, 248}
 
\subsubsection{AvR / AvC File Format Specification}
\noindent AvR files store a \textit{real-valued} averaged spectrum; AvC files store a \textit{complex-valued} averaged spectrum.
The file layout is identical except for the data section, which is $N \times 8$ bytes for AvR and $N \times 16$ bytes (interleaved Re/Im pairs) for AvC,
where $N$ is given by \texttt{fLength} in the TSingleArray descriptor.
Type conventions: fields prefixed \texttt{f} are Single-precision floats (4 bytes); all other 4-byte TSingleArray fields are U32.
 
\vspace{0.5em}
 
\begin{longtable}{%
  >{\bfseries}p{3.8cm}
  p{1.6cm}
  p{3.0cm}
  p{5.6cm}
}
 
\toprule
\rowcolor{headerblue}
\color{white}\normalfont\bfseries Field &
\color{white}\bfseries Bytes &
\color{white}\bfseries Type &
\color{white}\bfseries Description / Possible Values \\
\midrule
\endfirsthead
 
\toprule
\rowcolor{headerblue}
\color{white}\normalfont\bfseries Field &
\color{white}\bfseries Bytes &
\color{white}\bfseries Type &
\color{white}\bfseries Description / Possible Values \\
\midrule
\endhead
 
\bottomrule
\endfoot
 
%% Header
\multicolumn{4}{l}{\itshape\small\textbf{Header}} \\[2pt]
\rowcolor{sectiongray}
DataType        & 1  & TDataType & Response quantity (see enum below). \\
\rowcolor{rowlight}
HzRes           & 8  & Double   & Frequency resolution $\Delta f$ (Hz\,/\,bin). $> 0$. \\
StartFrequency  & 8  & Double   & Start frequency (Hz). $\geq 0$. \\
\rowcolor{rowlight}
EndFrequency    & 8  & Double   & End frequency (Hz). $>$ StartFrequency. \\
ScaleFactor     & 8  & Double   & Scale factor applied to the averaged data. \\
\rowcolor{rowlight}
nAverages       & 4  & U32      & Number of averages used. Typically 1--1000. \\
AverageType     & 1  & Enum     & Averaging method (see enum below). \\[6pt]
  
%% AverageType enum
\multicolumn{4}{l}{\itshape\small\textbf{AverageType} \textnormal{enum values}} \\[2pt]
\rowcolor{rowlight}
\quad RMS       & --- & value 0 & Root-mean-square averaging. \\
\quad Mean      & --- & value 1 & Linear (arithmetic mean) averaging. \\
\rowcolor{rowlight}
\quad Complex   & --- & value 2 & Complex (vector) averaging. \\
\quad Geometric & --- & value 3 & Geometric mean averaging. \\
\rowcolor{rowlight}
\quad None      & --- & value 4 & No averaging (single measurement). \\[6pt]
 
%% TSingleArray descriptor
\multicolumn{4}{l}{\itshape\small\textbf{TSingleArray} \textnormal{(MtxVec array descriptor --- all fields 4 bytes)}} \\[2pt]
\rowcolor{sectiongray}
fComplex        & 4  & Single   & \texttt{0.0} = real (AvR),\ \texttt{1.0} = complex (AvC). \\
\rowcolor{rowlight}
fLength         & 4  & Single   & Number of data points $N$ (stored as float). Determines data section size. \\
a{[2]}          & 4  & U32      & Reserved / alignment padding. \\
\rowcolor{rowlight}
fConditionCheck & 4  & Single   & Internal integrity / validation value. \\
Precision       & 4  & U32      & \texttt{0} = Single (32-bit), \texttt{1} = Double (64-bit). \\
\rowcolor{rowlight}
MtxVecVersion   & 4  & U32      & Version of MtxVec library that wrote the file. \\
SizeOf(TSample) & 4  & U32      & Byte size of each data sample. \\
\rowcolor{rowlight}
Tag             & 4  & U32      & User-defined tag or identifier. \\
MtxVecFileCode  & 4  & U32      & Magic number identifying this as a MtxVec file. \\
\rowcolor{rowlight}
a{[9]}          & 4  & U32      & Reserved / alignment padding. \\
a{[10]}         & 4  & U32      & Reserved / alignment padding. \\
\rowcolor{sectiongray}
caption         & 4  & String   & Short label, 4 bytes, null-padded (\texttt{\textbackslash 00\textbackslash 00\textbackslash 00\textbackslash 00} if empty). \\[6pt]
 
%% Data section
\multicolumn{4}{l}{\itshape\small\textbf{Data} \textnormal{(length determined by \texttt{fLength} and \texttt{fComplex})}} \\[2pt]
\rowcolor{rowlight}
AvR data        & $N \times 8$  & Double{[]}  & $N$ real-valued doubles, little-endian. Present when \texttt{fComplex = 0}. \\
\rowcolor{sectiongray}
AvC data        & $N \times 16$ & Double{[]}  & $N$ complex samples as interleaved Re\,/\,Im double pairs ($2N$ doubles total), little-endian. Present when \texttt{fComplex = 1}. \\[6pt]
 
%% Notes / comments
\multicolumn{4}{l}{\itshape\small\textbf{Notes / Comments} \textnormal{(appended after data section)}} \\[2pt]
\rowcolor{rowlight}
notes           & var & String   & Free-text annotation. Starts at byte offset $N \times 8$ (AvR) or $N \times 16$ (AvC) from the beginning of the data section, i.e.\ after \texttt{fLength} or $2 \times$\texttt{fLength} doubles respectively. \\
 
\end{longtable}
 
\vspace{1em}
\noindent\small\textbf{Data section size summary:}
\[
\text{AvR:}\quad N \times 8 \text{ bytes} \qquad \text{AvC:}\quad N \times 16 \text{ bytes} \qquad \text{where } N = \texttt{int(fLength)}
\]
 
\vspace{0.5em}
\noindent\small\textit{%
Notes: (1) \texttt{f}-prefixed TSingleArray fields are Single-precision floats (4 bytes).
(2) \texttt{MtxVecFileCode} and \texttt{SizeOf(TSample)} are U32.
(3) \texttt{caption} is 4 bytes, null-padded.
(4) All multi-byte values are little-endian.
(5) \texttt{TDataType} is stored as a single byte (U8).
Source: \texttt{george.pas} (Dew Research\,/\,DSP Master) and LabVIEW writer.}


\subsubsection{FRF File Format Specifications}
\noindent Fields verified against the Delphi source (\texttt{george.pas}) and LabVIEW writer code.
Type conventions: fields prefixed \texttt{f} are Single-precision floats (4 bytes);
\texttt{MtxVecFileCode} and \texttt{SizeOf(TSample)} are U32; \texttt{caption} is a null-padded string.
 
\vspace{0.5em}
 
\begin{longtable}{%
  >{\bfseries}p{3.8cm}   % Field
  p{1.6cm}               % Bytes
  p{3.0cm}               % Type
  p{5.6cm}               % Description / Possible Values
}
 
\toprule
\rowcolor{headerblue}
\color{white}\normalfont\bfseries Field &
\color{white}\bfseries Bytes &
\color{white}\bfseries Type &
\color{white}\bfseries Description / Possible Values \\
\midrule
\endfirsthead
 
\toprule
\rowcolor{headerblue}
\color{white}\normalfont\bfseries Field &
\color{white}\bfseries Bytes &
\color{white}\bfseries Type &
\color{white}\bfseries Description / Possible Values \\
\midrule
\endhead
 
\bottomrule
\endfoot
 
%%  TFRFHeader packed record
\multicolumn{4}{l}{\itshape\small\textbf{TFRFHeader} \textnormal{(Delphi packed record --- no padding between fields)}} \\[2pt]
\rowcolor{sectiongray}
index           & 4  & Integer  & Channel or measurement index. $\geq 0$. \\
\rowcolor{rowlight}
Xr              & 8  & Double   & Reference (input) channel number or physical X location. \\
Yr              & 8  & Double   & Response (output) channel number or physical Y location. \\
\rowcolor{rowlight}
XActual         & 8  & Double   & Physical X-axis value (e.g.\ excitation location). \\
YActual         & 8  & Double   & Physical Y-axis value (e.g.\ response location). \\
\rowcolor{rowlight}
FRFType         & 1  & TFRFType & \texttt{0} = \texttt{ftTransfer} (\texttt{.trf}),\ \texttt{1} = \texttt{ftPoint} (\texttt{.pnt}). \\
DataType        & 1  & TDataType & Response quantity (see enum below). \\
\rowcolor{rowlight}
HzRes           & 8  & Double   & Frequency resolution $\Delta f$ (Hz\,/\,bin). $> 0$. \\
StartFrequency  & 8  & Double   & Start frequency (Hz). $\geq 0$. \\
\rowcolor{sectiongray}
EndFrequency    & 8  & Double   & End frequency (Hz). $>$ StartFrequency. \\[6pt]
 
%%  TDataType enum
\multicolumn{4}{l}{\itshape\small\textbf{TDataType} \textnormal{enum values (U8)}} \\[2pt]
\rowcolor{rowlight}
\quad dtAccel   & --- & value 0 & Accelerance (acceleration\,/\,force). \\
\quad dtMob     & --- & value 1 & Mobility (velocity\,/\,force). \\
\rowcolor{rowlight}
\quad dtRecept  & --- & value 2 & Receptance (displacement\,/\,force). \\
\quad dtMic     & --- & value 3 & Microphone response. \\[6pt]
 
%%  TSingleArray — typed per field
\multicolumn{4}{l}{\itshape\small\textbf{TSingleArray} \textnormal{(MtxVec internal array descriptor --- all fields 4 bytes)}} \\[2pt]
\rowcolor{sectiongray}
fComplex        & 4  & Single   & Array-level complex flag: \texttt{0.0} = real, \texttt{1.0} = complex. \\
\rowcolor{rowlight}
fLength         & 4  & Single   & Number of elements in the data array (stored as float). \\
a{[2]}          & 4  & U32      & Reserved / alignment padding. \\
\rowcolor{rowlight}
fConditionCheck & 4  & Single   & Internal integrity / validation value. \\
Precision       & 4  & U32      & \texttt{0} = Single (32-bit), \texttt{1} = Double (64-bit). \\
\rowcolor{rowlight}
MtxVecVersion   & 4  & U32      & Version of MtxVec library that wrote the file. \\
SizeOf(TSample) & 4  & U32      & Byte size of each data sample. \\
\rowcolor{rowlight}
Tag             & 4  & U32      & User-defined tag or identifier. \\
MtxVecFileCode  & 4  & U32      & Magic number identifying this as a MtxVec file. \\
\rowcolor{rowlight}
a{[9]}          & 4  & U32      & Reserved / alignment padding. \\
a{[10]}         & 4  & U32      & Reserved / alignment padding. \\[6pt]
 
%%  Caption + FRF data
\multicolumn{4}{l}{\itshape\small\textbf{Payload}} \\[2pt]
\rowcolor{sectiongray}
caption         & 4  & String   & Short label, 4 bytes, null-padded (\texttt{\textbackslash 00\textbackslash 00\textbackslash 00\textbackslash 00} if empty). \\
\rowcolor{rowlight}
FRF data        & $N{\times}16$ & Double[] & Complex FRF samples as interleaved Re\,/\,Im DBL pairs, little-endian. $N$ = number of frequency lines. \\
 
\end{longtable}







\section {Adiabatic Wave Equation}
If the fluid is adiabatic, then
$$p=K\rho^\gamma$$
so $${\partial p\over \partial x}=K\gamma \rho^{\gamma-1}{\partial \rho\over \partial x}=c^2{\partial \rho\over \partial x}$$
and $${\partial p\over \partial t}=K\gamma \rho^{\gamma-1}{\partial \rho\over \partial t}=c^2{\partial \rho\over \partial t}$$
But from continuity
$$ {\partial  \rho \over \partial t} + \rho_0{\partial  u \over \partial x}=0$$
so
$$ {\partial  p \over \partial t} + {\rho_0 \over c^2}{\partial  u \over \partial x}=0$$
$$ {\partial  p \over \partial t} =- {\rho_0 \over c^2}{\partial^2  y \over \partial x\partial t}$$
or integrating by time,
$$ p  =- {\rho_0 \over c^2}{\partial  y \over \partial x} =- {\kappa}{\partial  y \over \partial x}$$
where $\kappa$ is the adiabatic bulk modulus


and
$$\rho_0 {\partial  u \over \partial t}+{\partial p\over \partial x} = 0$$
or
$${\partial p\over \partial x} = -\rho_0 {\partial^2  y \over \partial t^2}$$

and $$p=K\rho^\gamma$$
so $${\partial p\over \partial x}=K\gamma \rho^{\gamma-1}{\partial \rho\over \partial x}=c^2{\partial \rho\over \partial x}$$


\subsection{wave equation}
Assume that the pressure is $p$ and the speed of a fluid particle is $u$ with a mean of zero and the density everywhere is $\rho_0 + \rho$.  Then the 1D continuity gives us:
$$ {\partial (\rho_0 + \rho) \over \partial t} + {\partial (\rho_0 + \rho)u \over \partial x}=0$$
or
$$ {\partial  \rho \over \partial t} + \rho_0{\partial  u \over \partial x}+ {\partial (\rho u) \over \partial x}=0$$
and assuming the last term is small compared to the second to last, you get:
$$ {\partial  \rho \over \partial t} + \rho_0{\partial  u \over \partial x}=0$$
or
$$ {\partial^2  \rho \over \partial t^2} + \rho_0{\partial^2  u \over \partial x\partial t}=0$$
The momentum equation (1D, no gravity or friction terms) is:
$$ {\partial  ((\rho_0 + \rho)u) \over \partial t}+(\rho_0 + \rho)u{\partial u\over \partial x} +{\partial p\over \partial x} = 0$$
or
$$\rho_0 {\partial  u \over \partial t}+ \rho {\partial u \over \partial t}+ u {\partial \rho \over \partial t}+\rho_0 u{\partial u\over \partial x} + \rho u{\partial u\over \partial x}+{\partial p\over \partial x} = 0$$
Assuming the fluctuations are small:
$$\rho_0 {\partial  u \over \partial t}+\rho_0 u{\partial u\over \partial x} +{\partial p\over \partial x} = 0$$
Differentiate with respect to x:
$$\rho_0 {\partial^2  u \over \partial t \partial x}+\rho_0 u{\partial^2 u\over \partial x^2} +{\partial^2 p\over \partial x^2} = 0$$
Dropping the middle term since derivatives in $u$ should be about the same in both $x$ and $t$ and therefore the first term should be a lot bigger than the second - and combining both equations you get:
$$ {\partial^2 p\over \partial x^2} = {\partial^2  \rho \over \partial t^2} $$
If we assume an adiabatic process, then
$$p=K\rho^\gamma$$
So
$${\partial p\over\partial t}=\gamma K\rho^{\gamma-1}{\partial \rho\over \partial t}$$
or
$${\partial^2 p\over\partial t^2}=\gamma K\rho^{\gamma-1}{\partial^2 \rho\over \partial t^2}+ \gamma(\gamma-1) K\rho^{\gamma-2}({\partial \rho\over \partial t})^2$$
Linearizing loses the last term:
$${\partial^2 p\over\partial t^2}=\gamma K\rho^{\gamma-1}{\partial^2 \rho\over \partial t^2}$$
or 
$${\partial^2 p\over\partial t^2}=\gamma {p\over \rho^\gamma}\rho^{\gamma-1}{\partial^2 \rho\over \partial t^2}$$
or
$${\partial^2 p\over\partial t^2}=\gamma RT{\partial^2 \rho\over \partial t^2}$$
Combining then gives
$$ {\partial^2 p\over \partial x^2} = {\partial^2 p\over\partial t^2}{1\over \gamma RT} $$
or
$$ {\partial^2 p\over\partial t^2}={ \gamma RT} {\partial^2 p\over \partial x^2} $$

\subsection{Acoustic Impedance}
Next, we want to look at the acoustic impedance of the plate (acoustic impedance, $Z$, is the ratio of pressure to volume flow, specific acoustic impedance, $z$, is the ratio of pressure to volume flow per unit area - or acoustic flow velocity). 
$$p=ZU$$
For an infinitely long pipe, cross-sectional area S, we get:
$$Z={p\over U}={F\over SU}=...={\rho c\over S}$$ - I do not see it


Defining specific acoustic impedance to be:
$$p=zu$$
where a change in pressure ($p$) drives the fluid ($u$) like a voltage drives a current.

 then
$${{F/ A}\over u}=z$$
We can look at the forces on the plate, there should be the force due to pressure, a spring-like restoring force and a damping force due to the fact that the material is hardly a perfect spring:
$$\sum F = ma= pA - kx -\beta u$$
or
$$z={ma+ kx +\beta u\over Au} = {m\over A} {\ddot x\over \dot x} +{ \beta\over A} + {k\over A}{x\over \dot x}$$
If we assume a solution of the form
$$u=\dot x=Ce^{i\omega}$$
then the impedance becomes:
$$z= {m\over A} {i\omega Ce^{i\omega}\over Ce^{i\omega}} +{ \beta\over A} + {k\over A}{Ce^{i\omega}/ i\omega\over Ce^{i\omega}}$$
$$z= {m\over A} i\omega +{ \beta\over A} + {k\over A}{1\over i\omega}$$
or
$$z={ \beta\over A} + i\omega{m\over A}  - {i\over  \omega {A\over k}}$$
or 
$$z=r + i\omega M  - {i\over  \omega C}$$
where
$$ r = { \beta\over A}$$
$$M = {m\over A}$$
$$C={A\over k}$$

\subsection{Sound Intensity Measurements}
Sound intensity is defined as the sound power per unit surface area.  For a point source:
$$I={P\over 4 \pi r^2}$$
So the Sound Power Level is referenced to $P_0=10^{-12} W$ and the Sound Intensity Level is referenced to $I_0=10^{-12} W/m^2$
$$ SWL=10log_{10}(P/P_0) and SIL = 10log_{10}(I/I_0)$$
Since 
$$I=pu={p^2\over z}$$
then the sound pressure level is:
$$SPL=10log_{10}(p^2/p^2_0)=20log_{10}(p/p_0)$$


Sound hole is modeled with a frictional resistance $R_h$:
$$m_{hole-air}\ddot x_h = A_h \Delta p - R_h \dot x_h$$
and the top plate has a restoring force:
$$m_{plate}\ddot x = F - k_px_p -R_p\dot x_p + A_p \Delta p$$

\subsection{Quality factor}
Usually defined as 
$$Q = {f_{center}\over \Delta f} = {\sqrt mk \over R}= {\omega_0\over 2r}$$
were the spread of the peak ($\Delta f$) is typically measured by the 1/2 power reduction point, $r$ is the exponential decay of the signal and R is the damping coefficient ($m\ddot x + R\dot x + kx=0$)

\end{document}